\summarychapter{32}{Aiming at the Foundations}{The Cult of Reason}

\begin{inbrief}
After a century of religious wars something deep in Europe revolted. The Age of Reason (1648--1789)
put reason in the place of faith and happiness in this world in the place of preparation for the
next; it was the birth of secularism. The chapter traces its roots in the Renaissance, the wars of
religion and the new science, its moderate first generation (Locke) and its radical second (the
deists, Voltaire and the \emph{philosophes}), and the Christian reply.
\question{How did the spirit of the Age of Reason take root, and what did its attack on the
foundations of Christianity leave behind?}
\end{inbrief}

If Americans believed in saints, Benjamin Franklin would be one: practical, earthy, affable, witty
and above all tolerant. A few weeks before he died he told President Ezra Stiles of Yale that he had
``some doubts'' about the divinity of Jesus but had never studied the question and expected soon to
know the truth ``with less trouble''; he saw no harm in the belief if it made Jesus' teachings better
respected and observed. That is the spirit of the age. Dogma seemed hardly worth fretting about;
what mattered was behavior, tolerance and respect for those who differ. These attitudes live on in
the values of the Western world. The Reformation had shown again that faith and power are a potent
brew. As long as Christians held power they used it to compel conformity, and tens of thousands died.

\ssection{The Spirit of the Age of Reason}
The revolt, also called the Enlightenment, was an intellectual revolution, a new way of looking at
God, the world and oneself. In Charles Williams's words, national interests and mental relaxations
combined ``to exclude metaphysics from culture.'' The Middle Ages and the Reformation were ages of
faith in the sense that reason served faith and the mind obeyed authority, the church's for Catholics
and the Bible's for Protestants. God's Word came before human thought, and the chief concern of this
life was preparation for the next. The Age of Reason reversed this. The human mind, not faith,
emotion, myth or superstition, is the best guide to happiness here and now. Baron von Holbach stated
the program: scatter the clouds of ignorance, teach man to respect his own reason, to base his morals
on his own nature and the advantage of society, and to seek his happiness by promoting that of
others, so that he becomes ``a virtuous and rational being, who cannot fail to become happy.''

\ssection{The Renaissance}
The seeds lay in the Reformation era itself. \emph{Renaissance} means rebirth, the recovery of the
values of classical Greece and Rome in literature, politics and the arts. Its humanism sought human
flourishing, the well-rounded person, and held that reading the ancient texts, best of all in the
original languages, was a conversation with author and culture that builds the soul. The Reformation
is unthinkable without it, for every major Reformer came to believe that he must meet the Bible in
its original tongues. Its best representative was Erasmus of Rotterdam (1467--1536), whose
best-selling satires, above all \emph{Praise of Folly}, mocked monasticism and Scholasticism with
irony and common sense. His admirers said he laid the egg that Luther hatched. But in 1524 his
\emph{Diatribe on Free Will} exposed the gulf between them. Luther held that the will is enslaved and
cannot love or serve God without enabling grace; Erasmus thought this a dangerous doctrine that
relieves man of moral responsibility. The difference between Reformation and Renaissance lies there,
in the view of man. The Reformers preached original sin and a fallen world. The Renaissance thought
well of human nature and the universe, and that confidence flowered in the Enlightenment.

A second root was the century of religious conflict (1550--1650): the English Civil War, the
persecution of the Huguenots, the Thirty Years' War. Common decency rebelled against fanatical
clerics and against the burning or drowning of an old woman accused of witchcraft or heresy.
Religious prejudice came to seem a greater danger than atheism, and a thirst spread for tolerance and
for truths common to all.

\begin{sidebox}{Enlightenment}
The Enlightenment was really a series of enlightenments, reaching England, France and other lands
before it began in Germany. Its basic insight belongs to the first French thinkers, who were
fascinated by England. There they saw more prosperity and more liberty, and institutions of church
and state that restricted individuals less and left them freer to reason for themselves. They drew a
central conviction: people freed from the burden of cultural tradition and allowed to reason for
themselves will prosper. Freedom to think independently of tradition yields truth and prosperity, an
insight since extended in endless varieties.
\end{sidebox}

The third root was a new faith in law and order, born of modern science. Copernicus (1473--1543) put
the sun, not the earth, at the center. Kepler (1571--1630) concluded that a force from the sun moves
the planets in their courses. Galileo (1564--1642) examined the planets through a telescope and
proved that falling bodies accelerate at a constant rate. Isaac Newton (1642--1727) united these
discoveries. His \emph{Mathematical Principles of Natural Philosophy} (1687) brought all motion, in
heaven and on earth, under one master principle, the law of gravitation, and showed the universe as a
great machine running by unalterable laws.

\ssection{The Wonder of the World-Machine}
Europe's readers were captivated. The medieval world of angels and demons came to seem implausible
or superstitious beside a sun-centered universe explained by mathematics. Some felt that man had
been displaced from the center of creation, and God with him, since a world that runs by discernible
laws seems to need no sustainer. Shelley reminds the reader that modern science depended on Christian
convictions and was carried on by minds trained in Christian institutions, and that in these debates
God was reimagined as one more actor on the stage instead of the Lord of the whole. The vast universe
dismayed some. Pascal (1623--1662) recoiled from ``the terrifying expanses which engulf us as a
fleeting atom.'' Others rejoiced with Alexander Pope: ``God said, `Let Newton be!' and all was
light.'' If the universe is a machine with one grand design, man has only to think clearly to find
meaning and happiness. That confidence in the senses and reason gave the age its name. Not every
blacksmith or village priest turned intellectual; many lived and died in the faith of their fathers.
But the direction of Europe had changed.

For twelve hundred years Augustine's ideas had ruled: man is an enslaved sinner in need of
supernatural grace, and God ordained the state to protect truth and punish error. Now intellectuals
said that man is no sinner but a reasonable creature who needs common sense more than grace.
Christians met two climates. At the end of the seventeenth century, especially in England, some tried
to harmonize reason and faith: Christianity is wholly reasonable, with some truths, such as God's
existence, reached by reason and others, such as Christ's resurrection, by revelation. After 1700,
especially in France, confidence in reason soared and appeals to Scripture were dismissed as
superstition.

John Locke (1632--1704) best represents the first climate. His \emph{Essay Concerning Human
Understanding} calls God's existence ``the most obvious truth that reason discovers,'' yet his God
has little in common with the God of the Exodus or the resurrection. Attempts to justify Christianity
rationally, Shelley observes, often end by revising it to fit what seems rational at the time.
Christians do not believe as the conclusion of a long argument but because they have received the
testimony of the first eyewitnesses, vindicated by God's continuing power; reason may serve
revelation but cannot replace it. To his credit Locke accepted revelation, with fulfilled prophecy
and miracles as proof of Jesus' authority. But he found few dogmas necessary, all of them simple and
plain to ordinary people. Christianity has one essential doctrine, that Jesus is the Messiah, and one
demand, a good life, which Jesus reinforced with rewards and punishments, as reason itself requires.
Most traditional theology he set aside as irrelevant.

\ssection{The Rise of the Radicals}
The next generation felt less obligation to the Christian past and simply set revelation aside. The
parts of the Bible that agree with reason are unnecessary; the parts that contradict it, myths,
miracles and priestly mumbo jumbo, are untrue. Eighteenth-century Paris became the capital of a
cosmopolitan culture in which the leaders of Europe and the American colonies shared one community of
thought. There the \emph{philosophes} brought the age to its climax. They were not academic
philosophers but men of letters who analyzed society's evils and urged reforms, aiming to spread
knowledge and free the human spirit.

Atheism was not fashionable among them. Most of the famous ``infidels'' believed in a supreme being
but thought it superstition to say that he interferes with the world-machine. This is deism,
especially popular among English speakers, a halfway house on the road to atheism. Its God is the
watchmaker who made the world, wound it up and let it run. Being a perfect workman, he has no need to
intervene, so miracles and special revelation are rejected. Deists held that theirs was mankind's
original religion and that all others are distortions of it, invented by priests to enlarge their
own power.

The most influential propagandist was Voltaire (1694--1778), the embodiment of French skepticism. In
histories, plays, pamphlets, essays, novels and some ten thousand letters he popularized Newton's
science, fought for personal liberty and a free press, and spread the cult of reason. He was the most
relentless critic of the established churches, Protestant and Catholic, sickened by their intolerance
and petty quarrels. Yet he did not aim to destroy religion: if God did not exist, he said, it would
be necessary to invent him. His only rival was the French \emph{Encyclopedia} in seventeen volumes,
edited by Denis Diderot (1713--1784), the chief monument of the \emph{philosophes}. It proclaimed the
new science, championed tolerance, denounced superstition and expounded deism; its article on
Christianity professed respect for the religion of Jesus while leaving the reader with contempt for
its social failures. Unlike earlier critics, who were heretics or dissenters attacking the church in
Christ's name, these men attacked from outside, and aimed not at one dogma but at the foundation of
all Christian truth.

\ssection{Christianity on Trial}
They called Christianity a plot to hand the earth to a priestly caste, a scheme to exploit the
ignorant. Voltaire named it ``the infamous thing'' and pointed to the thousands of victims of its
intolerance. They judged it by simple human standards of good and evil: a church that sanctioned the
slaughter of fellow Christians for purity of doctrine was a wicked institution and an obstacle to
peace and progress. Shelley notes the irony that these ``human'' standards were themselves a legacy
of Christian teaching. Their weapon was truth: teach men, said Diderot, to hold as true only what
they have verified and proved. That standard excluded Christian doctrine from the start, since
arguments from the authority of Bible or church were not counted as reasonable. Miracles fared no
better. Once in the supernatural, Diderot argued, there are no limits: one man feeds five thousand
with five loaves, the next with one, and a third with the wind. The critics knew they were making a
revolution in Europe's beliefs, and Voltaire reported each triumph like a general.

\ssection{The Futility of Empty Confidence}
In Catholic countries such as France the reply was weak. Church leaders asked the authorities to
censor dangerous books and seldom understood the questions raised. In England several writers
answered deism effectively, above all Bishop Joseph Butler (1692--1752), whose \emph{Analogy of
Religion} virtually ended the debate for thinking people. The deists assumed they could read God's
wisdom and purpose plainly in nature. Butler did not argue for God's existence, which they granted,
nor reject reason, man's natural light; he challenged its sovereignty. Reason gives no complete
system, and in ordinary life it offers only probabilities. Nature is full of obscurities and
unexplained mysteries. If we meet difficulties there, we should not be surprised to meet them in
religion; and since we trust and act on the probable knowledge that experience gives us of nature,
it is reasonable to do the same in religion.

In the end deism fell by its own weakness. Its optimism was false: it had no explanation for evil
and disaster, and when asked why people do not always see the religious truths in nature it could
only blame the lies of priests, which convinced few. Its fall did not restore Christianity to the
center of Western culture. The negative work of the age endured. Art, education and politics were
freed from formal Christian influence in a deliberate attempt at a religiously neutral civilization,
with faith confined to church, home and heart, as in secular societies today. Christians are left
with a basic problem. How far may believers go, as citizens, in getting the state to enforce
Christian standards of conduct? And if they give that up, what norm should they seek to make binding
on everyone?

\begin{keyfacts}
\keypeople{Benjamin Franklin; Baron von Holbach; Erasmus; Copernicus; Kepler; Galileo; Isaac Newton;
Blaise Pascal; Alexander Pope; John Locke; Voltaire; Denis Diderot; Bishop Joseph Butler.}
\keydates{%
  \kd{1524}{Erasmus's \emph{Diatribe on Free Will} marks his break with Luther}
  \kd{1648--1789}{The Age of Reason}
  \kd{1687}{Newton's \emph{Mathematical Principles of Natural Philosophy}}
  \kd{1694--1778}{Voltaire}
  \kd{1704}{Death of John Locke}
  \kd{1752}{Death of Bishop Joseph Butler, author of \emph{The Analogy of Religion}}}
\keyterms{Age of Reason; Enlightenment; secularism; Renaissance humanism; world-machine; law of
gravitation; deism; watchmaker God; \emph{philosophes}; \emph{Encyclopedia}; ``the infamous thing'';
probability; religiously neutral civilization.}
\end{keyfacts}
